% Sección 1. Cifras: documentos/informe_resultados.md (secciones 1, 4.2 y 6), app/DATOS.md (inicio_rutas,
% inicio_salidas_*), Entrega1/README.md y experimentos/E7_elasticidad/README.md (sección 1).
\section{Introducción}
\label{sec:introduccion}

Sobusa opera buses urbanos en el área metropolitana de Barranquilla. El reto cubre tres de sus rutas,
SILENCIO, CRA~54 y CL~72 (62,9, 56,4 y 57,0~km de recorrido de catálogo), que comparten corredores con otras
empresas, con Transmetro y con el transporte informal. Hoy las salidas de cada ruta y hora se fijan con el
promedio de las ocho semanas anteriores para la misma ruta, hora y día de la semana, y el despacho se hace
desde el patio casi en tiempo real: un viaje se registra unos 12 minutos antes de salir. La regla arrastra el
pasado (un festivo contamina las ocho semanas siguientes) y no mira la carga de cada hora. Entre 2023 y 2025 la
demanda cayó de 15,31 a 12,87 millones de pasajeros, cerca de un 8~\% al año, con casi la misma flota. En 2025
la oferta coincidió con la que pide la carga del tramo más cargado en solo el 15,3~\% de las horas; sobraron
salidas en el 44,7~\% y faltaron en el 40,0~\%. Un día laboral de CRA~54 programa 288 salidas y su carga pide
282, pero 31 sobran en los valles y 29 faltan en los picos: el problema es la hora, no el total.

La pregunta del reto es cuántos buses despachar en cada ruta y hora, y cada cuánto. Tiene las tres etapas
clásicas de la planeación, la operación y el control del transporte en bus \citep{ibarra2015review}:
pronosticar los pasajeros de cada ruta, día y hora de salida con una semana de antelación, porque el horario se
arma cada semana; convertir el pronóstico en salidas por hora \citep{ceder1984frequency,ceder2016transit} que
equilibren operación, espera, hacinamiento y pasajes perdidos, y ejecutar el plan en el día \citep{daganzo2009bunching}.

\paragraph{Aporte.} Una cadena completa, de los datos a la decisión: (i) un pronóstico LightGBM de Poisson con la
serie de referencia corregida por fallas de registro y pesos para los datos corruptos, comparado con 51
candidatos, y un híbrido con el modelo fundacional Chronos-2; (ii) una medición causal de la elasticidad de los
pasajeros a las salidas, $\beta$, con un experimento natural; (iii) un optimizador entero de salidas por hora con
costos de fuente pública, evaluado con la demanda real de 2025; (iv) un simulador, un motor de despacho en tiempo
real y un modelo persona a persona del mismo día con nuestro plan; (v) el diseño de un piloto para medir $\beta$, y
(vi) una aplicación sin conexión, \comi{Sobusa Despacho}, que lee cada cifra de su archivo de resultados. Los
experimentos se citan por el código de su carpeta: E1--E13 para la demanda (E11 es el modelo final y E12 la
comparación de modelos), D1--D6 para el despacho (D5 es el optimizador y D6 el simulador) y P1--P2 para el
planificador y el piloto. Revisores independientes recalcularon sus cifras en tres rondas.
